\documentclass[10pt,a4paper]{amsart}
\usepackage[margin=19mm]{geometry}
\usepackage{amsmath,amssymb,mathtools}
\usepackage[scr=rsfs,cal=euler]{mathalfa}
\usepackage{xfrac}
\usepackage[unicode,hidelinks,bookmarks=false]{hyperref}
\newcommand{\ie}{i.e.\ }
\newcommand{\cf}{cf.\ }
\newcommand{\resp}{respectively\ }
\newcommand{\Subsection}{\subsection}
\providecommand{\nobreakdash}{\nobreak\hbox{-}}
\setlength{\emergencystretch}{2em}
\multlinegap0pt
\usepackage{bm}
\usepackage{bbm}
\usepackage{dsfont}
\usepackage{xspace}
\usepackage{cancel}
\usepackage{mathtools}
\usepackage{amsthm}
\makeatletter
\@ifundefined{theorem}{\newtheorem{theorem}{Theorem}}{}
\@ifundefined{lemma}{\newtheorem{lemma}[theorem]{Lemma}}{}
\@ifundefined{proposition}{\newtheorem{proposition}[theorem]{Proposition}}{}
\makeatother
\newcommand{\LC}{\mathrm{LC}}
\newcommand{\Rvm}{R_{\mathrm{vm}}}
\newcommand{\vm}{\leq_{\mathrm{vm}}}
\title[Vertex-minor Ramsey numbers: exact values and extremal struc]{Vertex-minor Ramsey numbers: exact values and extremal structure}
\author{Ji Ho Bae}
\date{Source: arXiv:2604.13434v1 (CC BY 4.0)}
\begin{document}
\maketitle

\noindent\emph{Source and licence.} Vertex-minor Ramsey numbers: exact values and extremal structure. Version arXiv:2604.13434v1 (CC BY 4.0), distributed under the Creative Commons Attribution 4.0 International licence, as recorded in the arXiv licence metadata for this exact version and shown on the abstract page https://arxiv.org/abs/2604.13434v1. Full rights notice: licenses/arxiv-2604.13434-CC-BY-4.0.txt.\par\smallskip

\noindent\textit{Editorial scope.} Counted unit: the Proposition whose source label is \texttt{prop:disconnected-k4} in the article (source0.tex lines 525--528, source label \texttt{prop:disconnected-k4}), delivered together with the article's own complete original proof of it (source0.tex lines 530--588, one \texttt{proof} environment of 59 lines). No mathematics is written, repaired or re-derived: statement and proof are the article's own text line for line. Required context retained uncounted: prop:alpha (lines 226--234); lem:disconnected (lines 264--279); prop:Rvm2 (lines 362--364); prop:Rvm3 (lines 388--390). Every definition, display and label of those lines is retained unchanged. Omitted from this package and not redistributed: 1--225, 235--263, 280--361, 365--387, 391--524, 529--529, 589--1068. Nothing inside a retained excerpt is omitted or altered: every retained body is delivered exactly as the article's lines are written, so no omission receipt of any kind accompanies this package. Citations: the retained excerpts carry no citation command at all, so no bibliography entry of the article is transcribed and no reference is redistributed. Reference adaptation: none was needed and none was made; every reference inside the delivered unit resolves inside it, so no unresolved reference or citation is produced. Printed slips kept literal: no printed word, symbol, hypothesis or number of the counted statement or of its proof was altered. Layout and class: the article's own class is \texttt{amsart} with packages including amsmath, amssymb, amsthm, booktabs, array, hyperref, xcolor; the wrapper is the lane's pinned \texttt{amsart} editorial wrapper at 10pt on A4, which restates the theorem-like environments the retained text needs and adds the article's own macro definitions that the retained text needs (\texttt{\textbackslash LC}, \texttt{\textbackslash Rvm}, \texttt{\textbackslash vm}), with extra wrapper packages (bm, bbm, dsfont, xspace, cancel, mathtools). The delivered numbering is the unit's own. Assets: no retained span contains a figure environment, a graphics-inclusion command, a PDF-page inclusion, a table or a TikZ picture; no image file is copied into this package, the package declares no asset dependency and no image file is redistributed with it. Licence: version arXiv:2604.13434v1 of this article is distributed under the Creative Commons Attribution 4.0 International licence, as recorded in the arXiv licence metadata for this exact version and shown on the abstract page https://arxiv.org/abs/2604.13434v1. A full rights notice is delivered at licenses/arxiv-2604.13434-CC-BY-4.0.txt. Characters: every retained excerpt is pure ASCII, so the delivered engine, XeLaTeX with its default Latin Modern text font, reports no missing character.
\begin{proposition}[LC-orbit characterization]\label{prop:alpha}
  Let $G$ be a graph and $k$ a positive integer.  Then
  \[
    E_k \vm G
    \quad\Longleftrightarrow\quad
    \alpha(G') \geq k
    \text{\; for some $G' \in [G]_{\LC}$.}
  \]
\end{proposition}

\begin{lemma}[LC orbit of a disconnected graph]\label{lem:disconnected}
  Let $G=G_1\cup\cdots\cup G_c$ be a graph whose connected components
  are $G_1,\dots,G_c$ on pairwise disjoint vertex sets.  Then
  \[
    [G]_{\LC}
    \;=\;
    \bigl\{\,G_1'\cup\cdots\cup G_c' :
             G_i'\in[G_i]_{\LC}\text{ for each }i\,\bigr\},
  \]
  and consequently
  $\;\displaystyle
    \max_{G'\in[G]_{\LC}}\alpha(G')
    \;=\;
    \sum_{i=1}^{c}\,\max_{G_i'\in[G_i]_{\LC}}\alpha(G_i').
  $
\end{lemma}

\begin{proposition}\label{prop:Rvm2}
  $\Rvm(2) = 3$.
\end{proposition}

\begin{proposition}\label{prop:Rvm3}
  $\Rvm(3) = 7$.
\end{proposition}

\begin{proposition}\label{prop:disconnected-k4}
  Every disconnected graph on at least $9$ vertices contains $E_4$
  as a vertex-minor.
\end{proposition}

\begin{proof}
  For a graph~$H$, write
  $\beta(H)=\max_{H'\in[H]_{\LC}}\alpha(H')$.  By
  Proposition~\ref{prop:alpha}, $E_4\vm H$ if and only if
  $\beta(H)\geq 4$.  From $\Rvm(2)=3$
  (Proposition~\ref{prop:Rvm2}) and $\Rvm(3)=7$
  (Proposition~\ref{prop:Rvm3}):
  \[
    \beta(H)\geq
    \begin{cases}
      3 & \text{if $|H|\geq 7$,}\\
      2 & \text{if $|H|\geq 3$,}\\
      1 & \text{if $|H|\geq 1$.}
    \end{cases}
  \]
  The third bound holds because $\alpha(H)\geq 1$ for every nonempty
  graph.

  Let $G=G_1\cup\cdots\cup G_c$ be a disconnected graph with
  $|G|=n\geq 9$ and $c\geq 2$ connected components, where
  $|G_i|=n_i$ and $n_1\geq\cdots\geq n_c\geq 1$.  By
  Lemma~\ref{lem:disconnected},
  $\beta(G)=\sum_{i=1}^{c}\beta(G_i)$.

  \smallskip
  \emph{Case~1: $n_1\geq 7$.}
  Then $\beta(G_1)\geq 3$.  Since $c\geq 2$, there exists $j\geq 2$
  with $\beta(G_j)\geq 1$, giving $\beta(G)\geq 3+1=4$.

  \smallskip
  \emph{Case~2: $n_i\leq 6$ for all~$i$.}
  We distinguish three subcases by the number of components of order
  at least~$3$.

  If at least two components satisfy $n_i\geq 3$, each contributes
  $\beta\geq 2$, so $\beta(G)\geq 2+2=4$.

  If exactly one component has $n_i\geq 3$, say~$G_1$ with
  $3\leq n_1\leq 6$, then every remaining component has $n_j\leq 2$.
  The total order of the remaining components is
  $n-n_1\geq 9-6=3$, so there are at least
  $\lceil 3/2\rceil=2$ of them, each contributing $\beta\geq 1$.
  Hence $\beta(G)\geq 2+2=4$.

  If every component has $n_i\leq 2$, then
  $c\geq\lceil n/2\rceil\geq 5$, and each contributes $\beta\geq 1$,
  so $\beta(G)\geq 5\geq 4$.

  \smallskip
  In every case $\beta(G)\geq 4$, so $E_4\vm G$ by
  Proposition~\ref{prop:alpha}.

  \smallskip
  The bound $n\geq 9$ is sharp: the graph
  $\overline{C_6}\cup K_2$ on $8$~vertices satisfies
  $\beta(\overline{C_6})=2$ (since $\Rvm(3)>6$) and
  $\beta(K_2)=1$, giving
  $\beta(\overline{C_6}\cup K_2)=3<4$.
\end{proof}

\end{document}
